\documentclass[12pt,a4paper]{article}
\usepackage{graphicx}
\usepackage{amsmath}
\usepackage{xcolor}
\usepackage{geometry}
\usepackage{fancyhdr}
\usepackage{tikz}
\usepackage{enumitem}

\geometry{a4paper, margin=1in}

% Set up the header and footer
\pagestyle{fancy}
\fancyhf{}
\fancyhead[L]{\leftmark}
\fancyhead[R]{\thepage}

\begin{document}

	% Title Page
	\begin{titlepage}
		\centering
		\vspace*{1cm}
		
		% Decorative Line
		\noindent\rule{\textwidth}{1pt}\\[0.3cm]
		
		% Project Title
		{\Huge\textbf{\color{blue}LAB NOTEBOOK USING LATEX}}\\
		\vspace{0.5cm}
		\noindent\rule{\textwidth}{1pt}\\[1.5cm]
		
		% Insert Image (adjusted size)
		\includegraphics[width=0.2\textwidth]{makaut.png} % Specify the path to your image here
		\vspace{0.8cm}
		
		% University Name
		\Large\textbf{\color{blue}Maulana Abul Kalam Azad University of Technology}\\
		\vspace{0.8cm}
		
		% Group Details
		\Large\textbf{Submitted by:}\\
		\vspace{0.3cm}
		
		\textbf{\color{red}Group Leader:}\\
		\vspace{0.2cm}
		\Large MD Erfanuddin\\
		\large Roll Number: 30001223001\\
		Department: BCA\\
		
		\vspace{1cm}
		
		\textbf{\color{green}Group Members:}\\
		\vspace{0.3cm}
		
		\large
		\begin{tabbing}
			\hspace{4cm} \= \hspace{8cm} \= \kill
			\textbf{Name} \> \textbf{Roll Number} \> \textbf{Department} \\
			\textcolor{purple}{Sanchita Laha} \> 30084323022 \> BSc in Data Science \\
			\textcolor{purple}{Arnov Pal} \> 30001223067 \> BCA \\
			\textcolor{purple}{Abir Dey} \> 30001223010 \> BCA \\
			\textcolor{purple}{Ibnoon Qais} \> 30059223003 \> BSc in Forensic Sciences \\
		\end{tabbing}
		
		\vfill
		
		% Department and University Information
		\Large\textbf{\color{orange}Git Assignment}\\
		
		% Decorative Line at the Bottom
		\vspace{1.5cm}
		\noindent\rule{\textwidth}{1pt}
		
	\end{titlepage}


	
	% Index Page
	\newpage
	\tableofcontents
	
	% Md Erfanuddin's Page
\newpage
\section{C Programming Calculator /by : Md. Erfanuddin}
\subsection{Introduction}
This document outlines how to create a C calculator program and push it to GitHub.

\subsection{Steps}
\begin{enumerate}
    \item Initialize local repository:
    \begin{verbatim}
        git init
    \end{verbatim}
    \item Build the C program.
    \item Commit the code:
    \begin{verbatim}
        git add .
        git commit -m "Initial commit: C calculator"
    \end{verbatim}
    \item Create a public repository on GitHub.
    \item Push to GitHub:
    \begin{verbatim}
        git remote add origin <repository-url>
        git push -u origin master
    \end{verbatim}
\end{enumerate}

\subsection{Code Snippet}
\begin{verbatim}
#include <stdio.h>

int main() {
    char operator; double num1, num2;
    printf("Enter operator (+, -, *, /): "); scanf("%c", &operator);
    printf("Enter two operands: "); scanf("%lf %lf", &num1, &num2);
    switch (operator) {
        case '+': printf("%.2lf + %.2lf = %.2lf\n", num1, num2, num1 + num2); break;
        case '-': printf("%.2lf - %.2lf = %.2lf\n", num1, num2, num1 - num2); break;
        case '*': printf("%.2lf * %.2lf = %.2lf\n", num1, num2, num1 * num2); break;
        case '/': if (num2 != 0.0) 
                      printf("%.2lf / %.2lf = %.2lf\n", num1, num2, num1 / num2);
                  else 
                      printf("Division by zero error.\n");
                  break;
        default: printf("Invalid operator.\n"); break;
    }
    return 0;
}
\end{verbatim}


% Abir Dey's Page
	\newpage
	\section{Enumerate ABC Format and Roman Number / by: Abir Dey}
	\subsection{Introduction}
	This document demonstrates how to create enumerated lists in LaTeX using ABC format and Roman numerals, providing flexibility in list formatting.
	
	\subsection{Creating an ABC Format Enumerated List}
	To create an enumerated list in ABC format (a, b, c), use the enumerate environment with label=\alph*.:
	
	\begin{verbatim}
		\begin{enumerate}[label=\alph*.]
			\item First item
			\item Second item
			\item Third item
		\end{enumerate}
	\end{verbatim}
	
	Result:
	
	\begin{enumerate}[label=\alph*.]
		\item First item
		\item Second item
		\item Third item
	\end{enumerate}
	
	\subsection{Creating a Roman Numerals Enumerated List}
	For a list with Roman numerals (i, ii, iii), use label=\roman*.:
	
	\begin{verbatim}
		\begin{enumerate}[label=\roman*.]
			\item First item
			\item Second item
			\item Third item
		\end{enumerate}
	\end{verbatim}
	
	Result:
	
	\begin{enumerate}[label=\roman*.]
		\item First item
		\item Second item
		\item Third item
	\end{enumerate}
% Arnov Pal's Page
 \newpage
 \section{Other Styles Such as Square and Circle /by: Arnov pal}
 \subsection{Introduction}
This document demonstrates the creation of geometric shapes like squares and circles using the TikZ package in LaTeX, commonly used in diagrams and illustrations.

\subsubsection{Square}
To create a square, we use the \texttt{rectangle} command in the \texttt{tikzpicture} environment, defining the bottom-left corner at (0,0) and the top-right corner at (3,3). This forms a square with a side length of 3 units, filled with light blue (\texttt{blue!30}) and outlined with a thick line.

\begin{verbatim}
\begin{tikzpicture}
 \draw[fill=blue!30, thick] (0,0) rectangle (3,3);
\end{tikzpicture}
\end{verbatim}

The result is shown below:

\begin{center}
 \begin{tikzpicture}
     \draw[fill=blue!30, thick] (0,0) rectangle (3,3);
 \end{tikzpicture}
\end{center}

\subsubsection{Circle}
For a circle, we use the \texttt{circle} command, specifying the center and radius. The code below generates a circle with a radius of 1.5 cm, filled with light red.

\begin{verbatim}
\begin{tikzpicture}
 \draw[fill=red!30, thick] (0,0) circle (1.5cm);
\end{tikzpicture}
\end{verbatim}

The result is shown below:

\begin{center}
 \begin{tikzpicture}
     \draw[fill=red!30, thick] (0,0) circle (1.5cm);
 \end{tikzpicture}
\end{center}

\subsection{Conclusion}
Creating and customizing shapes like squares and circles in LaTeX with TikZ allows for flexible and professional diagrams. Adjusting parameters like coordinates and colors can yield diverse results.


 % Sanchita Laha's Page
	\newpage
	\section{Create 5 Subsections/ by : Sanchita Laha}
	\subsection{Introduction}
	
	\subsubsection{Subsection 1: Overview}
	In this subsection, provide an overview of the topic you're discussing. This is the first of the five subsections.
	
	\subsubsection{Subsection 2: Background Information}
	Here, you can include any necessary background information that the reader needs to understand the topic.
	
	\subsubsection{Subsection 3: Methodology}
	Describe the methods or procedures used in your work or research. This is where you detail how you approached the problem.
	
	\subsubsection{Subsection 4: Analysis}
	Present the analysis of the data or information you've gathered. Discuss any patterns, trends, or insights you discovered.
	
	\subsubsection{Subsection 5: Conclusion}
	Summarize your findings and discuss the implications of your work. Mention any future directions or next steps.
 \vspace{0.5cm}
\section{How the Subsections were Created}
To create subsections in LaTeX, use the \texttt{\textbackslash subsection} command. Below is the LaTeX code used to create the subsections:

\begin{verbatim}
\subsection{Subsection 1: Overview}
Provide an overview of the topic you're discussing.

\subsection{Subsection 2: Background Information}
Include any necessary background information.

\subsection{Subsection 3: Methodology}
Describe the methods used in your research.

\subsection{Subsection 4: Analysis}
Present the analysis of the data or information.

\subsection{Subsection 5: Conclusion}
Summarize your findings and discuss next steps.
\end{verbatim}
    % Ibnoon Qais's Page
	\newpage
 \section{Mathematical Equations and Chemical Formulas /by: Ibnoon Qais}
 \subsection{Equation 1}
The equation for the square of a function is:
\begin{equation}
 f(x) = x^2  
\end{equation}
\textbf{LaTeX Code}: \texttt{f(x) = x\textasciicircum 2}

\subsection{Equation 2}
The equation in inline mode is:
\[
 f(x) = x^2
\]
\textbf{LaTeX Code}: \texttt{f(x) = x\textasciicircum 2}

\subsection{Equation 3}
The following are basic linear equations:
\begin{equation}
 1 + 2 = 3 \\
 x = 3 - 2
\end{equation}
\textbf{LaTeX Code}: 
\begin{verbatim}
\begin{equation}
 1 + 2 = 3 \\
 x = 3 - 2
\end{equation}
\end{verbatim}

\subsection{Equation 4}
The simplified form of the equations is:
\begin{equation*}
 1 = 3 - 2 \\
 x = 3 - 2
\end{equation*}
\textbf{LaTeX Code}:
\begin{verbatim}
\begin{equation*}
 1 = 3 - 2 \\
 x = 3 - 2
\end{equation*}
\end{verbatim}

\subsection{Equation 5}
Aligning multiple equations:
\begin{align*}
 1 + 2 &= 3 \\
 1 &= 3 - 2
\end{align*}
\textbf{LaTeX Code}: 
\begin{verbatim}
\begin{align*}
 1 + 2 &= 3 \\
 1 &= 3 - 2
\end{align*}
\end{verbatim}

\subsection{Equation 6}
Additional mathematical expressions and chemical formulas:
\begin{align*}
 f(x) &= x^2 \\                % Quadratic function
 g(x) &= \frac{1}{x} \\        % Reciprocal function
 F(x) &= \int_{b}^{a} \frac{1}{3}x^3 \, dx \\  % Definite integral
 H_{2}SO_{4} \\                % Chemical formula
 a^2 + b^2 &= c^2 \\           % Pythagorean theorem
 e^x &= \sum_{n=0}^{\infty} \frac{x^n}{n!}  % Exponential series
\end{align*}
\textbf{LaTeX Code}: 
\begin{verbatim}
\begin{align*}
 f(x) &= x^2 \\
 g(x) &= \frac{1}{x} \\
 F(x) &= \int_{b}^{a} \frac{1}{3}x^3 \, dx \\
 H_{2}SO_{4} \\
 a^2 + b^2 &= c^2 \\
 e^x &= \sum_{n=0}^{\infty} \frac{x^n}{n!}
\end{align*}
\end{verbatim}

    \end{document}
